\section{Introduction}
\label{sec:introduction}

A sensor calibration model is usually asked to do more than convert a voltage into a physical quantity. It must separate the quantity of interest from environmental influences, adapt to differences between instruments, and continue to work when some measurements disappear. These requirements are particularly visible in environmental monitoring, where inexpensive instruments are calibrated against a reference device for a limited period and then used without that reference~\citep{zimmerman2018calibration,devito2008airquality}. Calibration data are valuable, but they do not describe every combination of sensor availability that a deployment may encounter. A sensor may fail, be removed for servicing, or become unavailable at the same time as another sensor. A useful predictor therefore needs a distribution over the target for each available subset, together with a clear account of how the subset affects its uncertainty.

There are two different forms of adaptation in this problem. Across deployments, the response of a sensor can change in gain, offset, nonlinearity, and noise level. Within a deployment, the same collection of calibrated sensors can be observed in different subsets. Fitting a separate predictor for every subset uses the short calibration record repeatedly and does not scale well with the number of sensors. A single predictor with missing values filled by constants is easier to deploy, but it leaves the model to learn which parts of a filled input carry evidence and which parts merely encode absence. Neither choice directly explains how evidence from correlated measurements should be combined.

The correlation matters because several sensors can respond to a common disturbance. For example, an environmental variable can change several readings even when the target remains fixed. Once that disturbance has been marginalized, the readings are dependent conditional on the target. A product that assigns an independent likelihood factor to each sensor over the target alone then counts some shared evidence more than once. The resulting error is more specific than an overly narrow interval. Observing another sensor can change the optimal coefficient of an existing sensor, including its sign. A measurement that helps estimate the target in isolation can help estimate and subtract a disturbance when another measurement is present. Proposition~\ref{prop:scalar} makes this obstruction precise for fixed scalar Gaussian experts.

We consider prediction from a calibration support set. The predictor receives labeled calibration rows, including their observation masks, and a query containing the available sensor readings. Its parameters are trained across deployments, while its adaptation to a new deployment comes from that support set. This is the setting of in context probabilistic prediction studied by conditional neural processes and prior fitted networks~\citep{garnelo2018cnp,muller2022pfn,hollmann2025tabpfnv2}. Here the task distribution supplies repeated opportunities to learn calibration, but the prediction must still respect the measurement structure of the particular deployment.

Our construction represents each observed sensor by a Gaussian site over an augmented variable containing both the target and a small shared nuisance vector. The sites are multiplied with a Gaussian prior, and the target marginal is obtained by a small matrix solve. Under a linear Gaussian factor model, this is ordinary conditional Gaussian inference written in a form that omits an unavailable sensor by leaving out its site. Proposition~\ref{prop:exact} establishes exactness under the stated model and mask assumptions. Proposition~\ref{prop:rank} explains why the nuisance dimensions are needed for generic shared noise when the sites are fixed. Together, these results give a statistical reason for retaining shared residual variation in the fusion state.

The initial sites come from a factor analysis model fitted to the incomplete support. We then study two ways to learn corrections. The compact lifted cavity network encodes each sensor's calibration scatter and repeatedly updates its site using the posterior implied by the other sensors. The lifted cavity transformer uses a cell token transformer to read the support and query, then predicts corrections to the same Gaussian sites. Its backbone can capture interactions across the complete observed context, while the final aggregation retains a Gaussian representation with positive variance. Both models begin at the fitted factor analysis predictor because their correction heads are initialized at zero.

This separation between evidence representation and learned correction is useful for interpretation. Fixed anchor sites provide a conditional under a fitted statistical model. Learned corrections adapt the response slope, offset, nuisance loading, and noise variance to the current context. The resulting distribution retains a transparent fusion calculation even when those parameters are predicted by a neural network. The small model and transformer therefore test the same site representation with two different ways of extracting information from calibration examples.

The empirical study uses synthetic sensor deployments with known generating distributions and two real sensor networks. It compares learned refinement with statistical calibration models and transformers trained on the same synthetic task family. The compact model benefits from its structure on the main synthetic tasks, including queries with deletion counts absent from training and panels wider than the training panels. The transformer site layer tests the same representation with a more flexible context encoder and extends the evidence to another real target. Response shifts and real data comparisons help establish where each encoder is most useful. The evaluation preserves the distinction between original checkpoints and a later repair for factor sign ambiguity, and reports every endpoint in the committed protocols. Calibration quality, missing sensor behavior, and computational cost are evaluated as parts of the same prediction problem.
